% ============================================================================
% Informe Tarea 3 - Sistemas Distribuidos y Paralelos
% Universidad de Concepción
% Autores: Javier Torres, Juan Umaña, Benjamin Jimenez
% ============================================================================

\documentclass[12pt, a4paper]{article}

% --- Paquetes ---
\usepackage[utf8]{inputenc}
\usepackage[spanish,es-tabla]{babel}
\usepackage[margin=2.5cm]{geometry}
\usepackage{graphicx}
\usepackage{booktabs}
\usepackage{amsmath}
\usepackage{amssymb}
\usepackage{hyperref}
\usepackage{float}
\usepackage{pgfplots}
\pgfplotsset{compat=1.18}

\hypersetup{
    colorlinks=true,
    linkcolor=blue,
    citecolor=blue,
    urlcolor=blue
}

\begin{document}

% ============================================================================
% Encabezado
% ============================================================================
\begin{center}
    {\LARGE \textbf{Tarea 3: Procesamiento de Imágenes en CUDA}} \\[0.5cm]
    {\large Javier Torres \& Juan Umaña \& Benjamin Jimenez} \\[0.2cm]
    % TODO: Reemplazar con la URL real del repositorio.
    {\small Repositorio: \url{https://github.com/TU_REPOSITORIO_AQUI}} \\[0.5cm]
    {\large Sistemas Distribuidos y Paralelos} \\
    {\large Universidad de Concepción} \\
    {\large Junio 2026} \\[1cm]
\end{center}

% ============================================================================
% 1. Introducción
% ============================================================================
\section{Introducción}

El presente informe aborda el cálculo de la matriz de covarianza sobre el
dataset de imágenes DIV2K utilizando procesamiento en GPU mediante CUDA. La
matriz de covarianza es un componente esencial en técnicas de análisis
estadístico multivariado, como el Análisis de Componentes Principales (PCA),
y su cálculo eficiente en grandes volúmenes de datos constituye un desafío
tanto algorítmico como de gestión de recursos de hardware.

Dadas $N$ imágenes, cada una aplanada en un vector de dimensión $D$ (número
total de píxeles por canales de color), el procedimiento matemático se
descompone en las siguientes etapas:

\begin{enumerate}
    \item \textbf{Vector promedio.} Para cada componente $j \in \{1, \dots, D\}$
    del vector aplanado, se calcula la media aritmética sobre las $N$
    observaciones:
    \begin{equation}
        \mu_j = \frac{1}{N} \sum_{k=1}^{N} v_j^{(k)}
    \end{equation}
    donde $v_j^{(k)}$ denota la componente $j$-ésima de la imagen $k$.

    \item \textbf{Centrado de los datos.} Se sustrae el vector promedio a cada
    observación para obtener los datos centrados:
    \begin{equation}
        \overline{v}_j^{(k)} = v_j^{(k)} - \mu_j
    \end{equation}

    \item \textbf{Matriz de covarianza.} Se construye la matriz
    $\mathbf{C} \in \mathbb{R}^{D \times D}$ mediante la suma de productos
    punto de los vectores centrados:
    \begin{equation}
        C_{jj'} = \frac{1}{N-1} \sum_{k=1}^{N}
        \overline{v}_j^{(k)} \cdot \overline{v}_{j'}^{(k)}
    \end{equation}
\end{enumerate}

A medida que la resolución de las imágenes escala, la dimensión $D$ crece
rápidamente y con ella el tamaño de la matriz de covarianza ($D \times D$),
imponiendo restricciones severas sobre la memoria del dispositivo GPU (VRAM).
Este informe compara un enfoque tradicional sincrónico frente a una
orquestación asíncrona con CUDA Streams, evaluando su capacidad para superar
dichas limitaciones.

% ============================================================================
% 2. Descripción de Algoritmos
% ============================================================================
\section{Descripción de Algoritmos}

\subsection{Implementación Tradicional en CUDA}

En la implementación tradicional, toda la ejecución se canaliza a través del
Stream~0 (por defecto). Las transferencias de datos entre el host y el
dispositivo se realizan de forma sincrónica mediante \texttt{cudaMemcpy}, lo
que implica que cada operación de copia bloquea al host hasta su finalización
completa.

Para la multiplicación matricial subyacente al cálculo de la covarianza, se
emplea una estrategia de \textit{Tiling} con memoria compartida
(\textit{shared memory}). En este esquema, cada bloque de hilos carga
cooperativamente una submatriz (\textit{tile}) desde la memoria global hacia
la memoria compartida del multiprocesador de streaming (SM). Dado que la
memoria compartida posee una latencia significativamente menor que la memoria
global, cada elemento cargado puede ser reutilizado múltiples veces por los
hilos del bloque durante el cálculo del producto parcial, reduciendo el
número total de accesos a memoria global.

No obstante, este enfoque requiere que la totalidad de los datos ---la matriz
centrada $\overline{\mathbf{V}} \in \mathbb{R}^{N \times D}$ y la matriz de
covarianza resultante $\mathbf{C} \in \mathbb{R}^{D \times D}$--- resida
simultáneamente en la VRAM. Al escalar la resolución de las imágenes, la
demanda de memoria puede exceder rápidamente la capacidad física del
dispositivo, provocando errores de \textit{Out of Memory} (OOM).

\subsection{Orquestación con CUDA Streams}

Para superar la limitación de memoria y mejorar la utilización del hardware,
se diseñó una estrategia asíncrona basada en múltiples CUDA Streams. El
principio fundamental consiste en particionar los datos de entrada en lotes
(\textit{batches}) independientes, asignando cada lote a un stream distinto.

Los elementos clave de esta implementación son:

\begin{itemize}
    \item \textbf{Memoria paginada bloqueada (\textit{Pinned Memory}):}
    La memoria del host se asigna mediante \texttt{cudaMallocHost}, lo cual
    impide que el sistema operativo la intercambie a disco (\textit{swap}) y
    habilita el acceso directo a memoria (DMA) por parte del controlador de la
    GPU. Esta condición es prerrequisito para las transferencias asíncronas.

    \item \textbf{Transferencias asíncronas (\texttt{cudaMemcpyAsync}):}
    Las copias entre host y dispositivo se lanzan de forma no bloqueante dentro
    de cada stream. El host puede encolar operaciones en múltiples streams sin
    esperar la finalización de ninguna transferencia individual.

    \item \textbf{\textit{Latency hiding} (ocultamiento de latencia):}
    Gracias a la independencia entre streams, es posible solapar
    simultáneamente tres tipos de operaciones: la transferencia del lote $i+1$
    hacia la GPU, la ejecución del kernel sobre el lote $i$, y la copia de
    resultados del lote $i-1$ de vuelta al host. Este solapamiento oculta las
    latencias del bus PCIe tras el tiempo de cómputo.
\end{itemize}

De este modo, la VRAM solo necesita alojar los lotes activos en cada instante,
permitiendo procesar conjuntos de datos que exceden la memoria física del
dispositivo.

% ============================================================================
% 3. Metodología Experimental
% ============================================================================
\section{Metodología Experimental}

\subsection{Entorno de Ejecución}

Los experimentos se ejecutaron sobre un sistema operativo Linux con la
siguiente configuración de hardware:

\begin{table}[H]
    \centering
    \caption{Especificaciones del entorno de ejecución.}
    \label{tab:entorno}
    \begin{tabular}{ll}
        \toprule
        \textbf{Componente} & \textbf{Especificación} \\
        \midrule
        CPU          & % TODO: Completar (ej. Intel Core i7-10750H @ 2.60 GHz)
                       \\
        Memoria RAM  & % TODO: Completar (ej. 16 GB DDR4 2933 MHz)
                       \\
        GPU          & NVIDIA GeForce GTX 1650 Ti (3715 MB VRAM) \\
        Sistema Operativo & Linux % TODO: Distribución (ej. Ubuntu 22.04 LTS)
                            \\
        CUDA Toolkit & % TODO: Completar versión (ej. 12.2)
                       \\
        Compilador   & % TODO: Completar (ej. nvcc 12.2, g++ 11.4)
                       \\
        \bottomrule
    \end{tabular}
\end{table}

\subsection{Parámetros de Prueba}

Se evaluó el rendimiento del enfoque asíncrono variando el número de CUDA
Streams en el conjunto $S \in \{1, 2, 4, 8, 16\}$. Todas las pruebas se
realizaron utilizando imágenes escaladas a una resolución de $112 \times 84$
píxeles ---resolución que provoca un error de \textit{Out of Memory} bajo el
enfoque tradicional--- sobre la totalidad del dataset DIV2K.

% TODO: Indicar el número total de imágenes procesadas y la dimensión D
% resultante del aplanamiento.

% ============================================================================
% 4. Resultados
% ============================================================================
\section{Resultados}

\subsection{Tiempos de Ejecución}

El Cuadro~\ref{tab:tiempos} presenta los tiempos de ejecución desglosados
para cada configuración de streams evaluada.

\begin{table}[H]
    \centering
    \caption{Tiempos de ejecución (ms) para distintas configuraciones de
    streams ($112 \times 84$).}
    \label{tab:tiempos}
    \begin{tabular}{ccccc}
        \toprule
        \textbf{Streams ($S$)} &
        \textbf{Copia a GPU} &
        \textbf{Cómputo} &
        \textbf{Copia a Host} &
        \textbf{Tiempo Total} \\
        \midrule
        % TODO: Reemplazar los guiones con los tiempos medidos.
        1  & --- & --- & --- & --- \\
        2  & --- & --- & --- & --- \\
        4  & --- & --- & --- & --- \\
        8  & --- & --- & --- & --- \\
        16 & --- & --- & --- & --- \\
        \bottomrule
    \end{tabular}
\end{table}

\subsection{Speedup}

La Figura~\ref{fig:speedup} muestra el speedup obtenido en función del número
de streams, comparado con la línea de speedup ideal (lineal).

\begin{figure}[H]
    \centering
    \begin{tikzpicture}
        \begin{axis}[
            width=0.85\textwidth,
            height=0.55\textwidth,
            xlabel={Número de Streams ($S$)},
            ylabel={Speedup},
            xmode=log,
            log basis x=2,
            xtick={1, 2, 4, 8, 16},
            xticklabels={1, 2, 4, 8, 16},
            ymin=0,
            ymax=18,
            ytick={0, 2, 4, 6, 8, 10, 12, 14, 16},
            grid=major,
            grid style={dashed, gray!40},
            legend pos=north west,
            legend style={font=\small},
            tick label style={font=\small},
            label style={font=\small},
        ]

        % --- Línea de speedup ideal (punteada) ---
        \addplot[
            dashed,
            thick,
            color=gray,
            mark=none,
        ] coordinates {
            (1, 1) (2, 2) (4, 4) (8, 8) (16, 16)
        };
        \addlegendentry{Speedup ideal}

        % --- Speedup medido (línea sólida con marcadores) ---
        % TODO: Reemplazar las coordenadas Y con los speedups calculados
        %       (Speedup = T_1 / T_S, donde T_1 es el tiempo con 1 stream).
        \addplot[
            solid,
            thick,
            color=blue,
            mark=square*,
            mark size=3pt,
        ] coordinates {
            (1, 1.00)
            (2, 1.00)   % TODO: Reemplazar con speedup real
            (4, 1.00)   % TODO: Reemplazar con speedup real
            (8, 1.00)   % TODO: Reemplazar con speedup real
            (16, 1.00)  % TODO: Reemplazar con speedup real
        };
        \addlegendentry{Implementación}

        \end{axis}
    \end{tikzpicture}
    \caption{Speedup en función de la cantidad de CUDA Streams utilizados.}
    \label{fig:speedup}
\end{figure}

\subsection{Perfilado con Nsight Systems}

Para verificar visualmente el solapamiento entre transferencias y ejecución de
kernels, se utilizó el profiler NVIDIA Nsight Systems (\texttt{nsys}).

\begin{figure}[H]
    \centering
    % TODO: Reemplazar 'example-image' con la ruta a la captura de nsys.
    % Ejemplo: \includegraphics[width=\textwidth]{img/nsys_timeline.png}
    \includegraphics[width=\textwidth]{example-image}
    \caption{Diagrama temporal de Nsight Systems evidenciando el solapamiento
    entre transferencias Host$\leftrightarrow$GPU y ejecución de kernels con
    múltiples streams.}
    \label{fig:nsys}
\end{figure}

% ============================================================================
% 5. Discusión de Resultados y Análisis Crítico
% ============================================================================
\section{Discusión de Resultados y Análisis Crítico}

% --- Saturación del speedup ---
Los resultados del Cuadro~\ref{tab:tiempos} y la Figura~\ref{fig:speedup}
revelan que el speedup obtenido se aleja progresivamente de la línea ideal a
medida que se incrementa el número de streams. Este comportamiento se explica
por la naturaleza \textit{compute-bound} del algoritmo de covarianza: a
diferencia de problemas donde la transferencia de datos domina el tiempo total
(como algoritmos de ordenamiento \textit{memory-bound}), el cálculo del
producto $\mathbf{X}_c^{\top} \mathbf{X}_c$ concentra aproximadamente el 99\%
del tiempo de ejecución en operaciones aritméticas dentro de los núcleos CUDA.

Bajo estas condiciones, una vez que el solapamiento logra ocultar
completamente las latencias de copia ---lo cual se consigue con un número
relativamente bajo de streams---, incrementar la cantidad de streams no reduce
el componente de cómputo. La GPU se encuentra saturada ejecutando kernels, y
el speedup converge asintóticamente hacia un valor determinado por la razón
entre el tiempo de transferencia y el tiempo de cómputo.

% --- Mitigación del bus PCIe ---
En cuanto a la limitación del bus PCIe, la estrategia de \textit{latency
hiding} logró mitigarla de manera efectiva. Al solapar transferencias con
ejecución de kernels, el tiempo atribuible a las copias de datos queda oculto
tras el cómputo, eliminándolo como factor limitante del rendimiento global. Es
importante señalar que esta mitigación es viable precisamente porque el
algoritmo es intensivo en cómputo; en escenarios donde los tiempos de
transferencia y cómputo fueran comparables, el ancho de banda del bus PCIe
podría volver a constituir un cuello de botella significativo.

% --- Viabilidad out-of-core ---
El hallazgo más relevante de esta experimentación es la demostración de que la
orquestación con CUDA Streams constituye una técnica viable para el
procesamiento \textit{out-of-core}. El enfoque tradicional fracasa con un error
de \textit{Out of Memory} al intentar procesar imágenes de $112 \times 84$
píxeles en una GPU con 3715~MB de VRAM. En contraste, el enfoque asíncrono
---al particionar los datos en lotes y mantener en VRAM únicamente los
segmentos activos--- resuelve exitosamente el mismo problema sin comprometer la
corrección matemática del resultado. Esta capacidad de escalar más allá de los
límites físicos de la memoria del dispositivo resulta fundamental para
aplicaciones de procesamiento de imágenes a gran escala.

% ============================================================================
% 6. Conclusiones
% ============================================================================
\section{Conclusiones}

A partir del trabajo desarrollado, se extraen las siguientes conclusiones:

\begin{itemize}
    \item La implementación tradicional con Stream~0 y copias sincrónicas,
    si bien funcional para resoluciones bajas, resulta inviable al escalar la
    resolución de las imágenes debido a las limitaciones de VRAM del
    dispositivo.

    \item La orquestación con múltiples CUDA Streams, combinada con
    \textit{Pinned Memory} y transferencias asíncronas, permite superar las
    restricciones de memoria al procesar los datos en lotes, habilitando el
    cálculo de la covarianza para resoluciones que exceden la capacidad física
    de la VRAM.

    \item El speedup obtenido mediante \textit{latency hiding} es limitado
    por la naturaleza \textit{compute-bound} del algoritmo: una vez solapadas
    las latencias de transferencia, el rendimiento queda acotado por la
    capacidad de procesamiento de los núcleos CUDA.

    % TODO: Agregar conclusiones adicionales basadas en los resultados
    % numéricos obtenidos.
\end{itemize}

\end{document}
